% ============================================================
%  Docker Blockchain Registry — Présentation Beamer
%  Étudiants : EL-GHOUALI Aymane, ALMAMOUN Ouassama, DABBEBI Mohammed
%  Encadrant : BOITIER William
% ============================================================
\documentclass[aspectratio=169]{beamer}

% ---------- Thème & couleurs ----------
\usetheme{Madrid}
\usecolortheme{default}
\definecolor{dockerblue}{RGB}{13, 183, 237}
\definecolor{darkbg}{RGB}{15, 23, 42}
\definecolor{cardgray}{RGB}{30, 41, 59}
\definecolor{emerald}{RGB}{34, 197, 94}
\definecolor{amber}{RGB}{245, 158, 11}
\definecolor{red}{RGB}{239, 68, 68}

\setbeamercolor{structure}{fg=dockerblue}
\setbeamercolor{frametitle}{bg=darkbg, fg=white}
\setbeamercolor{title}{fg=white}
\setbeamercolor{subtitle}{fg=dockerblue}
\setbeamercolor{author}{fg=white}
\setbeamercolor{date}{fg=dockerblue}
\setbeamercolor{institute}{fg=white}
\setbeamercolor{titlelike}{bg=darkbg}
\setbeamercolor{palette primary}{bg=darkbg, fg=white}
\setbeamercolor{palette secondary}{bg=cardgray, fg=white}
\setbeamercolor{palette tertiary}{bg=dockerblue, fg=white}
\setbeamercolor{block title}{bg=dockerblue, fg=white}
\setbeamercolor{block body}{bg=cardgray!30, fg=black}
\setbeamercolor{block title alerted}{bg=red, fg=white}
\setbeamercolor{block body alerted}{bg=red!10, fg=black}
\setbeamercolor{block title example}{bg=emerald!80!black, fg=white}
\setbeamercolor{block body example}{bg=emerald!10, fg=black}

% ---------- Packages ----------
\usepackage[utf8]{inputenc}
\usepackage[T1]{fontenc}
\usepackage[french]{babel}
\usepackage{graphicx}
\usepackage{listings}
\usepackage{tikz}
\usepackage{fontawesome5}
\usepackage{array}
\usepackage{booktabs}
\usepackage{xcolor}

% ---------- Listings (code) ----------
\lstdefinestyle{solidity}{
  basicstyle=\tiny\ttfamily,
  keywordstyle=\color{dockerblue}\bfseries,
  commentstyle=\color{gray},
  stringstyle=\color{emerald},
  breaklines=true,
  frame=single,
  rulecolor=\color{dockerblue},
  backgroundcolor=\color{cardgray!15},
  morekeywords={function,mapping,address,bytes32,uint256,string,bool,
                event,modifier,require,emit,external,view,calldata,
                returns,pragma,solidity,contract,struct,public,private},
}

% ---------- Infos ----------
\title{\textbf{Docker Blockchain Registry}}
\subtitle{Registre décentralisé d'images Docker sur Ethereum}
\author{
  EL-GHOUALI Aymane \and
  ALMAMOUN Ouassama \and
  DABBEBI Mohammed
}
\institute{Encadrant : \textbf{BOITIER William}}
\date{Avril 2026}

% ============================================================
\begin{document}
% ============================================================

% ---- SLIDE 1 : TITRE ----------------------------------------
\begin{frame}
  \titlepage
  \vspace{-0.5cm}
  \begin{center}
    \small\textcolor{dockerblue}{\faDocker\ \ Smart Contract · Ethereum Sepolia · ECDSA Attestation}
  \end{center}
\end{frame}

% ---- SLIDE 2 : PROBLÈME ------------------------------------
\begin{frame}{Problème : Comment faire confiance à une image Docker ?}

  \begin{alertblock}{\faExclamationTriangle\ \ Risque réel dans les pipelines CI/CD}
    Une image Docker peut être \textbf{modifiée} entre sa construction et son déploiement
    sans que personne ne le détecte.
  \end{alertblock}

  \vspace{0.4cm}

  \begin{columns}[T]
    \column{0.48\textwidth}
      \textbf{\textcolor{red}{Sans solution :}}
      \begin{itemize}
        \item[\faTimesCircle] Image falsifiée en transit
        \item[\faTimesCircle] Registre Docker hackable
        \item[\faTimesCircle] Aucune preuve d'intégrité
        \item[\faTimesCircle] Déploiement de code malveillant
        \item[\faTimesCircle] Détection impossible a posteriori
      \end{itemize}

    \column{0.48\textwidth}
      \textbf{\textcolor{emerald}{Avec notre solution :}}
      \begin{itemize}
        \item[\faCheckCircle] Hash SHA256 ancré on-chain
        \item[\faCheckCircle] Vérification instantanée
        \item[\faCheckCircle] Falsification détectable
        \item[\faCheckCircle] Pipeline bloqué automatiquement
        \item[\faCheckCircle] Historique immuable
      \end{itemize}
  \end{columns}

  \vspace{0.4cm}
  \begin{block}{Attaque concrète}
    \small
    Un attaquant remplace \texttt{myapp:v1.0} dans le registre Docker par une image backdoorée.
    Sans vérification on-chain, le déploiement passe. \textbf{Notre contrat bloque cela.}
  \end{block}

\end{frame}

% ---- SLIDE 3 : CONCEPTS BLOCKCHAIN -------------------------
\begin{frame}{Concepts Blockchain essentiels}

  \begin{columns}[T]
    \column{0.48\textwidth}

      \begin{block}{\faLink\ \ Blockchain}
        Registre \textbf{distribué et immuable}. Les données sont
        répliquées sur des milliers de nœuds. Aucune entité centrale
        ne peut les modifier ou les supprimer.
      \end{block}

      \vspace{0.3cm}

      \begin{block}{\faFileCode\ \ Smart Contract}
        Programme autonome déployé on-chain. S'exécute \textbf{exactement
        comme codé}, sans intermédiaire. Logique et données résistent
        à toute censure.
      \end{block}

    \column{0.48\textwidth}

      \begin{block}{\faKey\ \ Cryptographie ECDSA}
        Signature à clé publique/privée. Permet de \textbf{prouver
        l'identité} d'un signataire sans révéler sa clé privée.
        Standard utilisé par Ethereum pour toutes les transactions.
      \end{block}

      \vspace{0.3cm}

      \begin{block}{\faHashtag\ \ Hash SHA256}
        Empreinte unique de 256 bits d'un fichier.
        \textbf{Toute modification} — même un seul bit — produit
        un hash complètement différent.
      \end{block}

  \end{columns}

\end{frame}

% ---- SLIDE 4 : ARCHITECTURE --------------------------------
\begin{frame}{Architecture de la solution}

  \begin{center}
  \begin{tikzpicture}[
    box/.style={draw, rounded corners, minimum width=2.8cm, minimum height=0.9cm,
                text centered, font=\small\bfseries},
    arrow/.style={->, thick, color=dockerblue},
    label/.style={font=\tiny, color=gray}
  ]
    % Nodes
    \node[box, fill=cardgray!40] (ci)      at (0,0)      {CI/CD Pipeline\\(GitHub Actions)};
    \node[box, fill=cardgray!40] (gateway) at (3.8,0)    {Gateway\\(Node.js)};
    \node[box, fill=dockerblue!20] (contract) at (7.6,0) {Smart Contract\\(Solidity)};
    \node[box, fill=cardgray!40] (ethereum) at (7.6,-2)  {Ethereum\\Sepolia};
    \node[box, fill=cardgray!40] (frontend) at (3.8,-2)  {Frontend\\(React + ethers.js)};
    \node[box, fill=cardgray!40] (user)    at (0,-2)     {Utilisateur\\(MetaMask)};

    % Arrows
    \draw[arrow] (ci)       -- node[above, label]{SHA256 + sign} (gateway);
    \draw[arrow] (gateway)  -- node[above, label]{registerImage()} (contract);
    \draw[arrow] (contract) -- node[right, label]{stockage} (ethereum);
    \draw[arrow] (user)     -- node[above, label]{vérification} (frontend);
    \draw[arrow] (frontend) -- node[above, label]{verifyImage()} (contract);
  \end{tikzpicture}
  \end{center}

  \vspace{0.2cm}
  \begin{columns}[T]
    \column{0.32\textwidth}
      \centering\textcolor{dockerblue}{\textbf{Enregistrement}}\\
      \small CI calcule SHA256,\\signe et soumet on-chain
    \column{0.32\textwidth}
      \centering\textcolor{emerald}{\textbf{Stockage}}\\
      \small Hash immuable sur\\Ethereum Sepolia
    \column{0.32\textwidth}
      \centering\textcolor{amber}{\textbf{Vérification}}\\
      \small Frontend compare\\hash local vs on-chain
  \end{columns}

\end{frame}

% ---- SLIDE 5 : SMART CONTRACT ------------------------------
\begin{frame}[fragile]{Smart Contract — DockerRegistry.sol}

  \begin{columns}[T]
    \column{0.52\textwidth}
      \begin{lstlisting}[style=solidity]
contract DockerRegistry {

  struct ImageRecord {
    bytes32 imageHash;
    uint256 timestamp;
    address registeredBy;
    bool    exists;
    bool    revoked;
    uint256 version;
    bool    signatureBased;
  }

  // N'importe qui peut enregistrer
  function registerImage(
    string  calldata imageName,
    bytes32          imageHash
  ) external { ... }

  // Registrant OU owner peut révoquer
  function revokeImage(
    string calldata imageName
  ) external { ... }

  // Vérification publique
  function verifyImage(
    string  calldata imageName,
    bytes32          imageHash
  ) external returns (bool) { ... }
}
      \end{lstlisting}

    \column{0.44\textwidth}
      \textbf{Fonctions clés :}
      \begin{itemize}
        \item \texttt{registerImage()} — ouvert à tous
        \item \texttt{registerImageWithSignature()} — ECDSA
        \item \texttt{revokeImage()} — registrant ou owner
        \item \texttt{verifyImage()} — public, gratuit
        \item \texttt{addSigner()} — gestion CI/CD
      \end{itemize}

      \vspace{0.3cm}
      \textbf{Événements on-chain :}
      \begin{itemize}
        \item \texttt{ImageRegistered}
        \item \texttt{ImageRevoked}
        \item \texttt{ImageVerified}
        \item \texttt{SignerAdded/Removed}
      \end{itemize}

      \vspace{0.3cm}
      \begin{exampleblock}{Déployé sur Sepolia}
        \tiny\texttt{0xc0Bb3dD54aebD5f86EE1B8\\eCC5c5d9B30743a9c0}
      \end{exampleblock}

  \end{columns}

\end{frame}

% ---- SLIDE 6 : ECDSA ATTESTATION --------------------------
\begin{frame}{Signature ECDSA — Attestation off-chain}

  \begin{block}{\faShieldAlt\ \ Pourquoi la signature ECDSA ?}
    Sans ECDSA, le CI doit avoir la \textbf{clé privée du contrat} dans ses secrets.
    Risque maximal si compromise. Avec ECDSA, le CI signe off-chain avec \textbf{sa propre clé}.
  \end{block}

  \vspace{0.3cm}

  \begin{columns}[T]
    \column{0.48\textwidth}
      \textbf{Flux de signature :}
      \begin{enumerate}
        \item CI calcule le SHA256 de l'image
        \item CI signe : \texttt{keccak256(name, hash,}\\\texttt{version, contractAddress)}
        \item N'importe qui soumet la signature on-chain
        \item Le contrat \textbf{récupère l'adresse} du signataire via \texttt{ECDSA.recover()}
        \item Vérifie que l'adresse est dans \texttt{trustedSigners}
      \end{enumerate}

    \column{0.48\textwidth}
      \textbf{Avantages :}
      \begin{itemize}
        \item[\faCheckCircle] Clé du contrat jamais exposée
        \item[\faCheckCircle] Multi-pipeline (GitHub, GitLab...)
        \item[\faCheckCircle] Révocation d'un signer sans redéployer
        \item[\faCheckCircle] Anti-replay : chaque signature unique
        \item[\faCheckCircle] Standard industrie (supply-chain)
      \end{itemize}

      \vspace{0.2cm}
      \begin{exampleblock}{Message signé}
        \tiny\texttt{keccak256(imageName \$\vert\vert\$ imageHash \$\vert\vert\$ version \$\vert\vert\$ contractAddr)}
      \end{exampleblock}
  \end{columns}

\end{frame}

% ---- SLIDE 7 : DÉCENTRALISATION ---------------------------
\begin{frame}{Décentralisation — Qui peut faire quoi ?}

  \begin{center}
  \renewcommand{\arraystretch}{1.4}
  \begin{tabular}{lccc}
    \toprule
    \textbf{Action} & \textbf{Tout utilisateur} & \textbf{Registrant} & \textbf{Owner} \\
    \midrule
    Enregistrer une image      & \textcolor{emerald}{\faCheck} & \textcolor{emerald}{\faCheck} & \textcolor{emerald}{\faCheck} \\
    Vérifier une image         & \textcolor{emerald}{\faCheck} & \textcolor{emerald}{\faCheck} & \textcolor{emerald}{\faCheck} \\
    Révoquer sa propre image   & \textcolor{red}{\faTimes}     & \textcolor{emerald}{\faCheck} & \textcolor{emerald}{\faCheck} \\
    Révoquer n'importe quelle  & \textcolor{red}{\faTimes}     & \textcolor{red}{\faTimes}     & \textcolor{emerald}{\faCheck} \\
    Ajouter un signataire CI   & \textcolor{red}{\faTimes}     & \textcolor{red}{\faTimes}     & \textcolor{emerald}{\faCheck} \\
    \bottomrule
  \end{tabular}
  \end{center}

  \vspace{0.4cm}

  \begin{columns}[T]
    \column{0.48\textwidth}
      \begin{exampleblock}{\faGlobe\ \ Données décentralisées}
        Stockées sur Ethereum — aucun serveur central,
        disponibles 24/7, impossibles à censurer ou supprimer.
      \end{exampleblock}

    \column{0.48\textwidth}
      \begin{block}{\faUsers\ \ Accès décentralisé}
        Tout utilisateur peut enregistrer et gérer
        ses propres images sans permission préalable.
      \end{block}
  \end{columns}

\end{frame}

% ---- SLIDE 8 : FRONTEND & CI/CD ---------------------------
\begin{frame}{Frontend \& Intégration CI/CD}

  \begin{columns}[T]
    \column{0.48\textwidth}
      \textbf{\faReact\ \ Frontend React (DApp)}
      \begin{itemize}
        \item React 18 + Vite + Tailwind CSS
        \item Connexion MetaMask (Sepolia)
        \item Dashboard : stats en temps réel
        \item Enregistrement / Vérification
        \item Gestion des signataires ECDSA
        \item Badges ECDSA vs Direct
        \item Switch réseau automatique
      \end{itemize}

    \column{0.48\textwidth}
      \textbf{\faGithub\ \ GitHub Actions}

      \textbf{register.yml} — au push sur main :
      \begin{enumerate}
        \item Build de l'image Docker
        \item Calcul SHA256
        \item Appel \texttt{gateway/push.js}
        \item Enregistrement on-chain
      \end{enumerate}

      \textbf{verify.yml} — sur Pull Request :
      \begin{enumerate}
        \item Rebuild de l'image
        \item Vérification on-chain
        \item \textcolor{red}{\textbf{Blocage}} si hash différent
      \end{enumerate}
  \end{columns}

  \vspace{0.3cm}
  \begin{alertblock}{\faExclamationCircle\ \ Pipeline bloqué automatiquement}
    Si le hash de l'image ne correspond pas à l'enregistrement blockchain,
    le déploiement est \textbf{bloqué avant} d'atteindre la production.
  \end{alertblock}

\end{frame}

% ---- SLIDE 9 : TESTS & RÉSULTATS --------------------------
\begin{frame}{Tests \& Résultats}

  \begin{columns}[T]
    \column{0.48\textwidth}
      \begin{exampleblock}{\faCheckDouble\ \ Suite de tests : 33/33 \faCheck}
        \begin{itemize}
          \item Déploiement \& initialisation
          \item Gestion des signataires (6 tests)
          \item \texttt{registerImage} direct (5 tests)
          \item \texttt{registerImageWithSignature} (8 tests)
          \item \texttt{revokeImage} (5 tests)
          \item \texttt{verifyImage} (5 tests)
        \end{itemize}
      \end{exampleblock}

    \column{0.48\textwidth}
      \textbf{Cas critiques testés :}
      \begin{itemize}
        \item[\faShieldAlt] \textbf{Anti-replay} : signature déjà utilisée rejetée
        \item[\faLock] \textbf{Signer révoqué} : refus immédiat
        \item[\faExclamationTriangle] \textbf{Hash falsifié} : \texttt{verifyImage} = false
        \item[\faUserSlash] \textbf{Tiers non autorisé} : révocation refusée
        \item[\faSync] \textbf{Re-registration} après révocation
      \end{itemize}

      \vspace{0.3cm}
      \begin{block}{Stack technique}
        Hardhat · ethers.js v6 · Chai · Mocha\\
        OpenZeppelin ECDSA · Solidity 0.8.24
      \end{block}
  \end{columns}

\end{frame}

% ---- SLIDE 10 : CONCLUSION ---------------------------------
\begin{frame}{Conclusion \& Perspectives}

  \begin{block}{\faTrophy\ \ Ce qui a été réalisé}
    \begin{itemize}
      \item Smart contract décentralisé avec \textbf{ECDSA signature attestation}
      \item Déployé sur \textbf{Ethereum Sepolia} — adresse permanente
      \item Frontend React complet avec connexion MetaMask
      \item Intégration CI/CD GitHub Actions (register + verify)
      \item 33 tests unitaires — couverture complète
    \end{itemize}
  \end{block}

  \vspace{0.3cm}

  \begin{columns}[T]
    \column{0.48\textwidth}
      \begin{exampleblock}{Problème résolu}
        Toute image Docker modifiée après enregistrement
        est \textbf{détectable instantanément} et bloque
        automatiquement le pipeline CI/CD.
      \end{exampleblock}

    \column{0.48\textwidth}
      \begin{block}{Perspectives}
        \begin{itemize}
          \item Déploiement frontend sur \textbf{IPFS}
          \item Support \textbf{multi-chaînes} (Polygon, Arbitrum)
          \item Intégration \textbf{Docker Hub} natif
          \item Vérification via \textbf{Sigstore/Cosign}
        \end{itemize}
      \end{block}
  \end{columns}

  \vspace{0.3cm}
  \begin{center}
    \textcolor{dockerblue}{\large\textbf{Merci pour votre attention}}\\
    \vspace{0.2cm}
    \small\textcolor{gray}{Code source disponible sur GitHub · Contrat :
    \texttt{0xc0Bb3dD54...743a9c0}}
  \end{center}

\end{frame}

% ============================================================
\end{document}
